% Side-by-side panels keep labels readable while fitting below the opening text.
\par\noindent\begin{minipage}{\linewidth}
\tagstructbegin{tag=Figure,alttext={左图 RNN 依次从春天到花再到开了传递状态；右图 Attention 从三个可见位置直接读取到位置三。右图线宽示意不同读取权重。}}
\begin{tcolorbox}[enhanced,colback=paper,colframe=linegray,
 boxrule=.7pt,arc=2mm,left=3mm,right=3mm,top=2mm,bottom=2mm]
\centering
\begin{tikzpicture}[>=Latex,x=1mm,y=1mm,
 word/.style={draw=linegray,fill=white,rounded corners=2pt,
  minimum width=17mm,minimum height=9mm,align=center,font=\footnotesize},
 output/.style={draw=flowcolor,fill=paper,rounded corners=2pt,
  minimum width=22mm,minimum height=10mm,align=center,font=\footnotesize}]
\node[font=\sffamily\bfseries\small,anchor=west] at (0,25) {RNN：依次传递};
\node[word] (r1) at (9,12) {春天\\$h_1$};
\node[word] (r2) at (30,12) {花\\$h_2$};
\node[word] (r3) at (51,12) {开了\\$h_3$};
\draw[->,very thick,flowcolor] (r1)--(r2);
\draw[->,very thick,flowcolor] (r2)--(r3);
\node[font=\footnotesize,align=center] at (30,-2) {较早信息经过中间状态};
\draw[linegray] (68,28)--(68,-9);
\node[font=\sffamily\bfseries\small,anchor=west] at (75,25) {Attention：直接读取};
\node[word] (a1) at (84,13) {春天\\$k_1,v_1$};
\node[word] (a2) at (108,13) {花\\$k_2,v_2$};
\node[word] (a3) at (132,13) {开了\\$k_3,v_3$};
\node[output] (at) at (108,-4) {位置 3 的\\读取结果};
\draw[->,flowcolor,line width=2pt] (a1.south) to[out=-80,in=165] (at.west);
\draw[->,flowcolor,line width=.6pt] (a2.south)--(at.north);
\draw[->,flowcolor,line width=1.15pt] (a3.south) to[out=-100,in=15] (at.east);
\end{tikzpicture}
\par
\end{tcolorbox}
\begin{center}\captionof{figure}{RNN 逐步传递与 Attention 直接读取的对照；线宽示意读取权重}\label{fig:attention-rnn}\end{center}
\tagstructend
\end{minipage}\par
